\section{Conclusiones}

% Todas las cifras salen de las tablas generadas por los scripts.

\begin{itemize}
  \item TCP obtuvo menor latencia que gRPC en las cuatro condiciones medidas.
  Con mensajes de 100~B en modo secuencial, la mediana del RTT de gRPC fue
  547,6~\% mayor que la de TCP (0,6171~ms frente a 0,0953~ms).
  \item La diferencia se redujo con mensajes grandes y con concurrencia: con
  10~KiB en modo concurrente, la mediana de gRPC fue solo 27,9~\% mayor, aunque
  su percentil~95 siguió siendo 145,4~\% mayor que el de TCP.
  \item El rendimiento de gRPC fue entre 55,9~\% y 89,6~\% menor que el de TCP,
  con la salvedad de que la duración de las tandas de gRPC incluye la llamada
  de control de cada envío.
  \item La hipótesis se cumplió en la parte de mensajes cortos; con 10~KiB la
  diferencia bajó, pero TCP mantuvo la menor mediana en todos los casos.
  \item En los 8\,000 mensajes de las 40 corridas medidas, el reloj de Lamport
  al recibir fue mayor que al enviar en el 100~\% de los casos, y el script
  encontró pares de eventos concurrentes con el mismo reloj, lo que confirma
  que Lamport no basta para detectar la concurrencia.
\end{itemize}
